% Incorporación aditiva para el Artículo I; no modifica la fuente anterior.
% RESULTADO_RECUPERADO: integral, capítulo 22, propietario extenso
% colaboracion/partes_i_ii/source/base_residencias/
% capitulo_22_c20_lineas_2613_3270.tex, líneas 176--576.
% Equivalencia 2/3: public/residencias/capitulo_22_voa_leech_moonshine.tex,
% líneas 55--127. Orden dos y Fricke: hmt/15d_voa_moonshine_orden_dos_rev11.tex.
% Prerrequisitos residentes en I: código C_W, matriz A_W, N(A_2^12),
% origen o_*, vector v_alpha de norma54, N_alpha reconocido como Leech,
% VOA reticular y presentación FLM de las secciones precedentes.
% Bibliografía nueva: chenlamshimakura2016, abelamyamada2017,
% kmconwaynorton1979 (bibitems listos para incorporar al final de este archivo).

\subsection{Orientation ternaire du voisin de Leech}
\label{km:orden-tres-reticular}

La présentation d’ordre deux construite précédemment utilise la négation
du réseau de Leech. Le même réseau, obtenu à partir de l’incidence marquée
par \(K\) et par le support dodécaphasique de \(\alpha\), admet une seconde
présentation compatible avec le caractère ternaire de la construction.
Pour expliciter cette continuité, on conserve le recollement, on construit
son isométrie d’ordre trois et on vérifie qu’elle préserve le voisin déjà
sélectionné. Cette étape n’exige de modifier ni \(K\), ni sa lecture rationnelle ni
la définition préalable de \(\alpha\). La proposition
\ref{km:estabilidad-vecino} vérifie l’isométrie réticulaire ; la proposition
\ref{km:traza-tipo} et le théorème~\ref{km:moonshine-tres} précisent la construction
d’ordre trois et les hypothèses des reconnaissances classiques utilisées.

Nous travaillons dans la base des racines simples de chaque facteur \(A_2\), avec
\begin{equation}
 B_2=\begin{pmatrix}2&-1\\-1&2\end{pmatrix},\qquad
 \mathsf C=\begin{pmatrix}0&-1\\1&-1\end{pmatrix},\qquad
 \mathsf R=\begin{pmatrix}0&1\\1&0\end{pmatrix},\qquad
 \rho=\begin{pmatrix}1\\1\end{pmatrix}.
 \label{km:coxeter-reflexion}
\end{equation}
La multiplication donne
\[
 \mathsf C^3=I_2,\quad
 \mathsf C^2+\mathsf C+I_2=0,\quad
 \mathsf C^{\mathsf T}B_2\mathsf C=B_2,\quad
 \mathsf R\mathsf C\mathsf R=\mathsf C^{-1},\quad
 \mathsf R^{\mathsf T}B_2\mathsf R=B_2.
\]
Ces identités fixent aussi bien la métrique que l'orientation de l'action. Dans le discriminant \(A_2^*/A_2\cong\mathbb F_3\), nous prenons pour générateur la classe de \(\varpi=(2,1)^{\mathsf T}/3\). On a \(\mathsf C\varpi-\varpi=(-1,0)^{\mathsf T}\), de sorte que \(\mathsf C\) agit trivialement sur le discriminant ; la réflexion \(\mathsf R\) change son signe.

La matrice de Paley--Witt déjà construite produit un élément de support
complet par l'opération suivante :
\begin{equation}
 \begin{gathered}
 q_*=(1,1,1,1,1,1),\qquad q_*A_W=(2,1,1,1,1,1),\\
 c_*=(q_*,q_*A_W)
 =(1,1,1,1,1,1,2,1,1,1,1,1)\in C_W.
 \end{gathered}
 \label{km:palabra-total}
\end{equation}
Chaque coordonnée de \(c_*\) est non nulle. Pour les douze positions,
nous définissons les cadres et l'isométrie
\begin{equation}
 P_b(c_*)=
 \begin{cases}\mathsf R,&(c_*)_b=1,\\I_2,&(c_*)_b=2,\end{cases}
 \qquad
 g_c=\bigoplus_{b=1}^{12}P_b(c_*)\mathsf C P_b(c_*)^{-1}.
 \label{km:isometria-ternaria}
\end{equation}
Le symbole \(P_b(c_*)\) désigne ici un repère inversible d’un facteur
\(A_2\), non le projecteur sur l’isotype tridimensionnel \(P_3\)
de la section~\ref{exc:k-direccion}.

\begin{proposition}[Isométrie d'ordre trois sur le voisin marqué]
\label{km:estabilidad-vecino}
L'opérateur \(g_c\) est une isométrie d'ordre trois sans vecteurs fixes non nuls. Il préserve le réseau de Niemeier \(N=N(C_W)\) et le voisin \(N_\alpha\) construit dans \eqref{exc:vecino}.
\end{proposition}
\begin{proof}
Chaque bloc de \(g_c\) est \(\mathsf C\) ou \(\mathsf C^{-1}\) ;
son polynôme minimal est \(X^2+X+1\). Cela prouve l'ordre et
l'absence de vecteurs fixes. Les deux blocs agissent trivialement sur le
discriminant, de sorte qu'ils conservent chaque classe de recollement et donc
le réseau \(N\).

Écrivons \(v=v_\alpha=4\rho_{o_*}+\sum_{b\ne o_*}\rho_b\),
avec les cadres et l'origine d'incidence déjà fixés. Ses coefficients
radiaux sont \(a_{o_*}=4\) et \(a_b=1\) pour \(b\ne o_*\).
Tous satisfont \(a_b\equiv1\pmod3\). Dans un facteur,
\[
 (\mathsf C-I_2)\rho=(-2,-1)^{\mathsf T},\qquad
 (\mathsf C^{-1}-I_2)\rho=(-1,-2)^{\mathsf T}.
\]
Après division par trois, leurs classes discriminantes sont respectivement \(2\) et \(1\). En conséquence, les vecteurs
\begin{equation}
 y=\frac{g_cv-v}{3},\qquad z=\frac{g_c^{-1}v-v}{3}
 \label{km:desplazamientos}
\end{equation}
ont pour mots discriminants \(c_*\) et \(-c_*\), tous deux dans
\(C_W\). Par conséquent, \(y,z\in N\). Le produit local est
\[
 \left\langle
 \frac{(\mathsf C^{\pm1}-I_2)a_b\rho}{3},a_b\rho
 \right\rangle=-a_b^2,
\]
et leur somme donne \(\langle y,v\rangle=\langle z,v\rangle=-(16+11)=-27\). Ainsi, \(y,z\in N_0=\{x\in N:\langle x,v\rangle\equiv0\pmod3\}\). Pour \(x\in N_0\),
\[
 \langle g_cx,v\rangle
 =\langle x,g_c^{-1}v\rangle
 =\langle x,v\rangle+3\langle x,z\rangle\equiv0\pmod3.
\]
L'argument appliqué à \(g_c^{-1}\) prouve \(g_cN_0=N_0\).
Enfin, \(g_c(v/3)=v/3+y\) conserve la classe qui engendre
\(N_\alpha=N_0+\mathbb Z(v/3)\). On obtient
\(g_cN_\alpha=N_\alpha\).
\end{proof}

La preuve conserve plus que la dimension du réseau : elle utilise le code de recollement, son orientation, l'origine du drapeau et la congruence du vecteur radial. Ces données expliquent pourquoi la seconde présentation reste liée au registre dodécaphasique dont la construction est partie.

\subsection{Trace ternaire, poids tordu et extension d'ordre trois}
\label{km:orbifold-tres}

Soit \(\Lambda=N_\alpha\) et soit \(\widehat g_c\) le relèvement
standard de \(g_c\) à l'algèbre de réseau
\(V_\Lambda=M(1)\otimes\mathbb C_\varepsilon[\Lambda]\), avec
l'extension centrale et le cocycle déjà déclarés. L'ordre impair permet
de choisir ce relèvement d'ordre trois. Chaque facteur complexe \(A_2\)
apporte les valeurs propres \(\zeta_3\) et \(\zeta_3^2\), où
\(\zeta_3^2+\zeta_3+1=0\).

\begin{proposition}[Trace graduée et type du relèvement]
\label{km:traza-tipo}
Pour \(j=1,2\),
\begin{equation}
 \operatorname{Tr}_{V_\Lambda}
       \bigl(\widehat g_c^{\,j}q^{L_0-1}\bigr)
 =t_3(\tau)
 :=q^{-1}\prod_{n\ge1}(1+q^n+q^{2n})^{-12}.
 \label{km:tres-traza}
\end{equation}
Le poids minimal des modules tordus non triviaux est \(4/3\), et
le relèvement est de type zéro.
\end{proposition}
\begin{proof}
Dans la partie réticulaire de la trace, seul le vecteur nul contribue,
car \(g_c\) et \(g_c^2\) n'ont pas de vecteurs fixes non nuls.
À chaque degré d'oscillateur, la trace des puissances symétriques est l'inverse
du déterminant \(I-g_c^jq^n\). Chaque facteur \(A_2\) contribue
\(1+q^n+q^{2n}\), ce qui démontre \eqref{km:tres-traza}.

La formule du poids du vide tordu pour un automorphisme de réseau d'ordre trois, appliquée aux deux espaces propres de dimension douze, donne
\[
 \rho_{\widehat g_c}
 =\frac{1}{4\cdot3^2}
   \bigl(1\cdot2\cdot12+2\cdot1\cdot12\bigr)=\frac43.
\]
Le type est la classe de \(3^2\rho_{\widehat g_c}=12\) modulo
trois, qui est nulle. La formule de poids est un résultat de la théorie des
modules tordus ; les multiplicités et l'isométrie qui l'alimentent
ont été construites ici.
\end{proof}

\begin{theorem}[Seconde présentation de l'algèbre Moonshine]
\label{km:moonshine-tres}
L'extension cyclique de type zéro
\begin{equation}
 V_{(3)}=\operatorname{Orb}_{\mathbb Z/3}
                 (V_\Lambda,\widehat g_c)
 \label{km:extension-orden3}
\end{equation}
est isomorphe à l'algèbre d'opérateurs de vertex \(V^\natural\).
\end{theorem}
\begin{proof}
Le théorème \ref{exc:teorema-leech} fournit un réseau positif,
pair, unimodulaire, de rang vingt-quatre et sans racines. La
Proposition~\ref{km:estabilidad-vecino} construit sur ce même
réseau une isométrie d'ordre trois sans vecteurs fixes ; la
Proposition~\ref{km:traza-tipo} vérifie le relèvement et son type.
On applique alors le théorème de Chen--Lam--Shimakura sur l'
orbifold d'ordre trois de l'algèbre du réseau de Leech
\cite[Teorema~1.1]{chenlamshimakura2016}. L'identification est également comprise
dans le traitement uniforme d'Abe--Lam--Yamada
\cite[Teorema~4.4 y secciones~2--4]{abelamyamada2017}. Cette application est postérieure à la construction
des données du réseau ; l'égalité des caractères ne remplace pas les
hypothèses du théorème utilisé.
\end{proof}

Pour préciser la symétrie duale, écrivons \(V_{(3)}=W_0\oplus W_1\oplus W_2\), où \(W_i\) est le secteur entier sélectionné dans le module \(\widehat g_c^{\,i}\)-tordu. Le produit satisfait \(Y(W_i,z)W_j\subset W_{i+j\bmod3}((z))\). Par conséquent,
\begin{equation}
 g^\vee|_{W_i}=\zeta_3^i I
 \label{km:dual-ciclico}
\end{equation}
définit un automorphisme d'ordre trois. En effet, le facteur sur un
produit est \(\zeta_3^{i+j}=\zeta_3^i\zeta_3^j\), et le vide et le
vecteur conforme appartiennent à \(W_0\). L'identification classique de
cet automorphisme dans le Monstre est la classe \(3B\), de série
\(T_{3B}=t_3+12\) \cite{kmconwaynorton1979,borcherds}.
Cette normalisation peut être vérifiée dans la construction. Le caractère
du réseau est \(\operatorname{ch}V_\Lambda=J+24\) : son pôle initial
est \(q^{-1}\), son espace de poids un est de dimension 24 et le
caractère modulaire de poids zéro est déterminé par ces données.
Si \(F\) est le caractère du sous-espace fixe et \(H\) celui de chaque
secteur tordu entier, la projection sur les invariants et la conjugaison
des deux secteurs donnent
\[
 F=\frac{J+24+2t_3}{3},\qquad F+2H=J,
 \qquad \operatorname{Tr}g^\vee q^{L_0-1}=F-H=t_3+12.
\]
Son domaine est l'algèbre étendue, tandis que \(\widehat g_c\) agit sur l'algèbre de réseau précédente.

\subsection{Équivalence des présentations d'ordres deux et trois}
\label{km:equivalencia-dos-tres}

Notons
\[
 V_{(2)}=(V_\Lambda)^+\oplus(V_\Lambda^T)^+
\]
la présentation FLM de \eqref{exc:flm}. Les deux présentations conservent le même voisin de Leech comme antécédent. Leurs extensions emploient des automorphismes différents et des produits de vertex déterminés au moyen des données de relèvement respectives.

\begin{theorem}[Équivalence graduée et transport d'automorphismes]
\label{km:equivalencia-voa}
Des isomorphismes d'algèbres d'opérateurs de vertex étant fixés
\[
 \phi_2:V_{(2)}\xrightarrow{\ \cong\ }V^\natural,
 \qquad
 \phi_3:V_{(3)}\xrightarrow{\ \cong\ }V^\natural,
\]
la composition
\begin{equation}
 \Phi_{23}=\phi_2^{-1}\phi_3:
 V_{(3)}\xrightarrow{\ \cong\ }V_{(2)}
 \label{km:Phi23}
\end{equation}
conserve le vide, le vecteur conforme, la graduation et le produit de vertex. La conjugaison par \(\Phi_{23}\) identifie les groupes d'automorphismes et conserve leurs traces graduées.
\end{theorem}
\begin{proof}
La composition des deux isomorphismes préserve les structures
déclarées. En particulier,
\[
 \Phi_{23}Y_{(3)}(a,z)b
 =Y_{(2)}(\Phi_{23}a,z)\Phi_{23}b,
 \qquad
 \Phi_{23}L_0^{(3)}=L_0^{(2)}\Phi_{23}.
\]
Si \(h\) est un automorphisme de \(V_{(3)}\), alors
\(h' =\Phi_{23}h\Phi_{23}^{-1}\) est un automorphisme de
\(V_{(2)}\). L'égalité des traces sur chaque composante homogène
de dimension finie donne
\[
 \operatorname{Tr}_{V_{(3)}}h q^{L_0-1}
 =\operatorname{Tr}_{V_{(2)}}h' q^{L_0-1}.
\]
\end{proof}

Le choix de \(\phi_2\) et \(\phi_3\) fait partie de l'énoncé ;
\(\Phi_{23}\) n'est pas présenté comme une identification canonique
indépendante de ces choix. Il ne transforme pas non plus un automorphisme
d'ordre trois en l'involution d'ordre deux : la conjugaison conserve
l'ordre. Son contenu est que deux extensions distinctes réalisent la
même structure graduée, avec un transport explicite de leurs produits,
représentations et traces.

\subsection{Fonctions êta, involution de Fricke et normalisation du vide}
\label{km:fricke}

Pour \(\operatorname{Im}\tau>0\), les traces des deux automorphismes de réseau admettent les expressions
\begin{equation}
 \eta(\tau)=q^{1/24}\prod_{n\ge1}(1-q^n),\qquad
 t_2=\left(\frac{\eta(\tau)}{\eta(2\tau)}\right)^{24},\qquad
 t_3=\left(\frac{\eta(\tau)}{\eta(3\tau)}\right)^{12}.
 \label{km:eta-trazas}
\end{equation}
Elles se déduisent directement des factorisations
\(1-q^{2n}=(1-q^n)(1+q^n)\) et
\(1-q^{3n}=(1-q^n)(1+q^n+q^{2n})\). L'ordre en \(q\) du
quotient de niveau trois, avant de l'élever à la puissance douze, est
\[
 \frac1{24}-\frac3{24}=-\frac1{12};
 \qquad 12\left(-\frac1{12}\right)=-1.
\]
Ici \(-1/12\) désigne la différence exacte des exposants
initiaux de deux fonctions êta. L’opération est spécifiée et ne
requiert pas d’interpréter une somme divergente comme une somme convergente.
L’égalité de ce scalaire avec un défaut normalisé d’une autre
réalisation n’identifie pas les opérateurs qui le produisent : chaque
réalisation conserve son domaine et son opération.

\begin{proposition}[Involution de Fricke sur la trace d'ordre deux]
\label{km:fricke-prop}
Soit \(w_2\tau=-1/(2\tau)\). Alors
\begin{equation}
 t_2(w_2\tau)=\frac{2^{12}}{t_2(\tau)},\qquad
 T_{2A}=t_2+24+\frac{2^{12}}{t_2}
 \label{km:fricke-t2}
\end{equation}
et \(T_{2A}(w_2\tau)=T_{2A}(\tau)\).
\end{proposition}
\begin{proof}
La loi de transformation
\(\eta(-1/\tau)=(-i\tau)^{1/2}\eta(\tau)\) donne
\[
 \frac{\eta(-1/(2\tau))}{\eta(-1/\tau)}
 =\sqrt2\,\frac{\eta(2\tau)}{\eta(\tau)}.
\]
La puissance vingt-quatre produit la première identité. La substitution
échange les deux termes variables de \(T_{2A}\) et laisse fixe
le terme constant. La reconnaissance de cette expression comme série
de McKay--Thompson de la classe \(2A\) correspond à la
symétrisation de niveau deux de Conway--Norton
\cite[\S\S4 y 6, tabla~3]{kmconwaynorton1979} et au théorème
du Moonshine \cite{borcherds}.
\end{proof}

Les uniformisations modulaires de niveaux deux et trois, dans les
normalisations précédentes, sont
\begin{align}
 j&=\frac{(t_2+256)^3}{t_2^2},
 \label{km:j-nivel2}\\
 J=j-744
 &=t_3+12+\frac{10\cdot3^9}{t_3}
       +\frac{4\cdot3^{14}}{t_3^2}
       +\frac{3^{18}}{t_3^3}.
 \label{km:j-nivel3}
\end{align}
Elles sont utilisées comme identités modulaires classiques sur les traces déjà
construites, et non comme règles pour sélectionner le code ou le voisin.
La seconde s'écrit également
\(j=(t_3+27)(t_3+243)^3/t_3^3\).
Le développement direct du produit de \eqref{km:tres-traza} commence par
\[
 t_3=q^{-1}-12+54q-76q^2-243q^3+1188q^4+O(q^5).
\]
En particulier, \([q]J=54+10\cdot3^9=196884\) ; c'est une conséquence des traces et de l'uniformisation. La construction de \(V^\natural\) conserve sa preuve par les réseaux et les orbifolds, indépendamment de cette vérification du premier coefficient.

% BEGIN R27_MG09
\subsubsection{Congruence ternaire de la trace et optimalité de l'exposant}
\label{mg09:congruencia}

L'identité de niveau trois permet de comparer tous les coefficients
de la trace construite à ceux de sa réalisation modulaire. Le domaine
approprié pour cette conséquence est l'anneau des séries formelles : chaque
coefficient dépend d'un nombre fini de facteurs du produit,
tandis que l'énoncé conserve tous les degrés.

\begin{lemma}[Unité entière du produit d'ordre trois]
La série
\begin{equation}
 u(q)=q\,t_3(q)=\prod_{n\geq1}(1+q^n+q^{2n})^{-12}
 \label{mg09:unidad}
\end{equation}
appartient à \(1+q\mathbb Z[[q]]\), et son inverse appartient au
même ensemble.
\end{lemma}

\begin{proof}
Chaque polynôme \(1+q^n+q^{2n}\) a pour terme constant un.
Pour toute série \(v(q)=1+\sum_{r\geq1}v_rq^r\) à
coefficients entiers, l'équation \(v(q)w(q)=1\) détermine
récursivement
\[
 w_0=1,\qquad w_r=-\sum_{i=1}^{r}v_iw_{r-i}\quad(r\geq1).
\]
Cette récurrence prouve l'existence, l'unicité et le caractère entier de l'inverse. Appliquée à chaque facteur, elle prouve également le caractère entier de sa puissance inverse douze. En outre, \((1+q^n+q^{2n})^{-12}\in1+q^n\mathbb Z[[q]]\). Pour calculer un coefficient de degré au plus \(r\), les facteurs avec \(n\leq r\) suffisent ; les produits partiels se stabilisent coefficient par coefficient. Ainsi, \eqref{mg09:unidad} définit une série à coefficients entiers et à terme constant un. La même récurrence, appliquée maintenant à \(u\), donne \(u^{-1}\in1+q\mathbb Z[[q]]\).
\end{proof}

\begin{proposition}[Divisibilité uniforme exacte]
Pour la trace \(t_3\) et la série \(J=j-744\) reliées par
\eqref{km:j-nivel3}, on a
\begin{equation}
 J-(t_3+12)\in3^9q\mathbb Z[[q]].
 \label{mg09:congruencia-exacta}
\end{equation}
L'exposant uniforme de divisibilité par trois est exactement neuf :
\begin{equation}
 [q]\bigl(J-t_3-12\bigr)=10\cdot3^9=196830.
 \label{mg09:optimalidad}
\end{equation}
\end{proposition}

\begin{proof}
Comme \(t_3=q^{-1}u\), le lemme donne
\(t_3^{-r}=q^ru^{-r}\in q^r\mathbb Z[[q]]\) pour
\(r=1,2,3\). L'identité \eqref{km:j-nivel3} se transforme en
\begin{equation}
 J-t_3-12
 =3^9q\left(10u^{-1}+4\cdot3^5q\,u^{-2}
                         +3^9q^2u^{-3}\right).
 \label{mg09:factorizacion}
\end{equation}
Tous les termes entre parenthèses sont des séries entières ; cela prouve
\eqref{mg09:congruencia-exacta}. Son terme constant est \(10\),
car \(u(0)=u^{-1}(0)=1\), et les deux derniers termes contiennent
\(q\) et \(q^2\). On obtient \eqref{mg09:optimalidad}. Comme
\(3\nmid10\), le coefficient de \(q\) a une valuation
\(3\)-adique de neuf, ce qui établit l'optimalité de l'exposant uniforme.
\end{proof}

En particulier, \([q^r]J\equiv[q^r]t_3\pmod{3^9}\) pour tout
\(r\geq1\). Le quantificateur procède de l'identité formelle et de
l'inversion entière démontrée dans le lemme. Une vérification jusqu'à
un degré fini quelconque est une spécialisation de cette relation.
La trace procède toujours de l'automorphisme d'ordre trois du réseau
construit à partir de l'incidence HMT ; l'uniformisation agit ensuite
sur cette trace. La congruence conserve cet ordre et détermine une
relation exacte entre ses deux lectures graduées.

% END R27_MG09
% BEGIN REV09_FRICKE_JET
\subsection{Correspondance de Fricke et récupération différentielle}
\label{corpus9:iv:fricke}

La trace réticulaire de niveau deux et sa complétion de Fricke
fournissent, après la construction exceptionnelle, les
coordonnées
\[
 t=t_2,\qquad s=\frac{4096}{t},\qquad
 T=t+s+24,\qquad t\in\mathbb C^\times.
\]
L’involution échange \(t\) et \(s\). Les deux lectures modulaires
ordonnées sont
\[
 j_1=\frac{(t+256)^3}{t^2},\qquad
 j_2=\frac{(s+256)^3}{s^2}.
\]
Ces formules utilisent l’uniformisation déjà établie dans
\eqref{km:j-nivel2} et conservent le choix de la lecture orientée.

\begin{proposition}[Correspondance et discriminant]
\label{corpus9:iv:fricke-correspondencia}
On a les identités
\begin{align}
 j_1-j_2&=(s-t)(T+23),\\
 j_1+j_2&=T^2+T-7256,\\
 j_1j_2&=(T+248)^3.
 \label{corpus9:iv:fricke-simetricos}
\end{align}
Les deux valeurs sont racines de
\[
 X^2-(T^2+T-7256)X+(T+248)^3=0,
\]
dont le discriminant est
\begin{equation}
 \Delta=(T-152)(T+104)(T+23)^2.
 \label{corpus9:iv:fricke-discriminante}
\end{equation}
\end{proposition}
\begin{proof}
Le développement de la première lecture, en utilisant \(ts=4096\), donne
\(j_1=t+768+48s+s^2\) ; la seconde s’obtient en échangeant
\(s,t\). La différence donne \((s-t)(t+s+47)\), et la somme se réduit à
\((t+s)^2+49(t+s)-6656\). La substitution de \(t+s=T-24\) produit
les deux premières identités. Pour le produit,
\[
 j_1j_2
 =\frac{((t+256)(s+256))^3}{(ts)^2}
 =(T+248)^3.
\]
Enfin,
\((s-t)^2=(T-24)^2-16384=(T-152)(T+104)\).
Multiplier par \((T+23)^2\) prouve
\eqref{corpus9:iv:fricke-discriminante}.
\end{proof}

\begin{proposition}[Récupération de l’orientation par les valeurs et le premier jet]
\label{corpus9:iv:fricke-jet}
Pour \(T\ne-23\), la paire ordonnée \((T,j_1)\) récupère
\begin{equation}
 t=\frac{T^2-j_1-3904}{T+23}.
 \label{corpus9:iv:fricke-inversa}
\end{equation}
Dans \(T=-23\), les deux paramètres

\[
 t_\pm=\frac{-47\pm45i\sqrt7}{2}
\]
partagent \((T,j_1)=(-23,-3375)\). Le premier jet
\(p=dj_1/dT\), calculé sur la branche paramétrée par \(t\),
les sépare et permet de récupérer
\begin{equation}
 p_\pm=\frac{-45\mp45i\sqrt7}{2},\qquad
 s=p-1,\qquad t=-46-p.
 \label{corpus9:iv:fricke-jet-inversa}
\end{equation}
\end{proposition}
\begin{proof}
De \(s^2=(T-24)s-4096\) résulte

\[
 j_1=(T+23)s+T-3352
     =T^2-3904-(T+23)t.
\]
Pour \(T\ne-23\), cette égalité donne
\eqref{corpus9:iv:fricke-inversa}. Pour \(T=-23\),
les conditions \(t+s=-47\), \(ts=4096\) donnent
\(t^2+47t+4096=0\), de discriminant \(-45^2\cdot7\).
La même identité donne \(j_1=-3375\) aux deux racines.


La dérivée \(dT/dt=1-4096/t^2\) ne s’annule qu’en
\(t=\pm64\). Aucune des deux racines exceptionnelles n’a cette
valeur, de sorte que \(T\) est une coordonnée locale sur les deux branches.
En différentiant la première expression de \(j_1\),
\[
 \frac{dj_1}{dt}
  =(s+1)\frac{dT}{dt}+(T+23)\frac{ds}{dt}.
\]
Dans \(T=-23\), il vient \(p=s+1=-46-t\), qui prouve
\eqref{corpus9:iv:fricke-jet-inversa}, avec des signes opposés
pour les indices appariés \(t_\pm,p_\pm\).
\end{proof}

Les points fixes de Fricke, \(t=\pm64\), expriment respectivement
\((T,j)=(152,8000)\) et \((-104,1728)\). La coïncidence en
\(T=-23\) correspond en revanche à deux points distincts de la
paramétrisation. Le corps de définition s’obtient à partir de l’équation :
\[
 \mathbb Q(t_\pm)=\mathbb Q(\sqrt{-7}),\qquad
 \left(\frac{2t+47}{45}\right)^2=-7.
\]
Sur cette fibre, Fricke coïncide avec la conjugaison quadratique,
de trace \(-47\) et de norme \(4096\). Le jet conserve ce corps,
car \((2p+45)^2=-45^2\cdot7\). La reconstruction est exacte
dans cette correspondance modulaire : la donnée différentielle distingue
les deux orientations qui partagent leurs valeurs scalaires.
% END REV09_FRICKE_JET

% BEGIN REV09_NIVELES_2_5
\subsection{Compatibilité des lecteurs modulaires de niveaux deux et cinq}
\label{corpus9:iv:niveles}

Sur la réalisation modulaire postérieure à la construction HMT,
fixons \(\operatorname{Im}\tau>0\), \(q=e^{2\pi i\tau}\) et
\(q^{1/5}=e^{2\pi i\tau/5}\). Le lecteur de Rogers--Ramanujan est

\[
 R(q)=\frac{q^{1/5}}{1+\dfrac{q}{1+\dfrac{q^2}{1+\cdots}}},
 \qquad x=R(q)^5.
\]
L’identité modulaire de cette réalisation, avec la normalisation
de \(j\) utilisée dans \eqref{km:j-nivel2}, est
\begin{equation}
 j=-\frac{(x^4-228x^3+494x^2+228x+1)^3}
          {x(x^2+11x-1)^5}.
 \label{corpus9:iv:nivel5-identidad}
\end{equation}
L’identité constitue la prémisse de réalisation modulaire du
transport algébrique suivant ; les germes, l’orientation et
les représentations d’auto-échelle restent ceux de APP--TRIT--TPK.

\begin{proposition}[Quotient pentadique et sélection de branche]
\label{corpus9:iv:nivel5-cociente}
Pour \(x\ne0\), soit \(u=x-x^{-1}\). Dans le domaine \(u\ne-11\),
\begin{equation}
 j=-\frac{(u^2-228u+496)^3}{(u+11)^5}.
 \label{corpus9:iv:nivel5-u}
\end{equation}
L’involution \(x\mapsto-1/x\) conserve \(u\). Ses préimages sont
\[
 x=\frac{u\pm\sqrt{u^2+4}}2.
\]
Le revêtement quadratique se ramifie en \(u=\pm2i\) ; l’application
rationnelle ultérieure \(u\mapsto j\) est de degré six.
\end{proposition}
\begin{proof}
Les factorisations exactes sont
\[
 \begin{aligned}
 x^4-228x^3+494x^2+228x+1
   &=x^2(u^2-228u+496),\\
 x^2+11x-1&=x(u+11).
 \end{aligned}
\]
Leur substitution dans \eqref{corpus9:iv:nivel5-identidad} prouve
\eqref{corpus9:iv:nivel5-u}. L’équation de la fibre est
\(x^2-ux-1=0\) ; son discriminant donne les points de ramification.
Le numérateur de \eqref{corpus9:iv:nivel5-u} est de degré six
et le dénominateur de degré cinq. En \(u=-11\), le polynôme
du numérateur vaut \(3125\), de sorte qu’ils n’ont pas de
facteur commun. Cela prouve le degré indiqué.
\end{proof}

Pour \(u\) réel, les deux racines ont des signes opposés. Dans le
domaine complexe, la branche choisie dans le paramétrage
par \(\tau\) est conservée. La lecture \(j\) seule conserve une fibre
génériquement de six valeurs de \(u\).


\begin{proposition}[Limite d’auto-échelle et compatibilité des niveaux]
\label{corpus9:iv:niveles-ramas}
Dans la limite réelle \(q\to1^-\),
\[
 R(q)\longrightarrow\varphi_{\mathrm{HMT}}^{-1},\qquad
 x\longrightarrow\varphi_{\mathrm{HMT}}^{-5}
     =\frac{5\sqrt5-11}{2},\qquad u\longrightarrow-11.
\]
Lorsque \(t=t_2(\tau)\) et \(u\) sont lus sur le même paramètre
modulaire, ils satisfont
\begin{equation}
 (t+256)^3(u+11)^5
   +t^2(u^2-228u+496)^3=0.
 \label{corpus9:iv:niveles-enlace}
\end{equation}
En posant \(\epsilon=u+11\), les branches qui satisfont les limites
indiquées ont les lois
\begin{align}
 \epsilon\to0,\ t\to0:
 &\quad \frac{t^2}{(-\epsilon)^5}\longrightarrow
          \frac{2^{24}}{5^{15}},\\
 \epsilon\to0,\ t\to\infty:
 &\quad t\epsilon^5\longrightarrow-5^{15}.
 \label{corpus9:iv:niveles-asintotica}
\end{align}
\end{proposition}
\begin{proof}
Pour \(0<q<1\), les réduites positives alternées de la
fraction continue encadrent \(R(q)\). Chaque réduite fixée
tend, lorsque \(q\to1^-\), vers la réduite correspondante de
la fraction continue dont tous les numérateurs sont égaux à un.
Les deux sous-suites de ces réduites ont pour limite
\(\varphi_{\mathrm{HMT}}^{-1}\), dans sa réalisation quadratique
déjà établie. L’encadrement prouve la première limite. Élever à
la puissance cinq et utiliser \(\varphi^2=\varphi+1\) donne les
deux autres.


Égaler \eqref{km:j-nivel2} et \eqref{corpus9:iv:nivel5-u}
prouve \eqref{corpus9:iv:niveles-enlace}. La substitution
\(u=\epsilon-11\) donne
\[
 \frac{(t+256)^3}{t^2}
   =-\frac{(3125-250\epsilon+\epsilon^2)^3}{\epsilon^5}.
\]
Dans la branche \(t\to0\), l’identité se réarrange sous la forme
\[
 \frac{t^2}{(-\epsilon)^5}
   =\frac{(t+256)^3}{(3125-250\epsilon+\epsilon^2)^3}.
\]
La limite est \(256^3/3125^3=2^{24}/5^{15}\).
Sur la branche \(t\to\infty\), multiplier par \(\epsilon^5\)
et écrire le membre de gauche comme
\(t\epsilon^5(1+256/t)^3\) donne la seconde limite.
\end{proof}

La relation conserve les branches d’une correspondance sur le
même \(j\). Le choix de \(\tau\), ou du relèvement
correspondant, fixe la paire transportée. Le pôle
pentadique et l’inversion de Fricke maintiennent leurs domaines et leurs
opérations respectifs.
% END REV09_NIVELES_2_5


Les trois opérations qui apparaissent dans ce prolongement ont des domaines distincts. La symétrie duale \(g^\vee\) agit sur des secteurs de l'algèbre de vertex ; Fricke agit sur le paramètre modulaire et sa fonction ; la dualité de rayon échange impulsion et enroulement dans le secteur réticulaire compact. Leurs formules de réciprocité peuvent être coordonnées au moyen des applications correspondantes, mais le nom « dualité » ne remplace pas ces applications. Cette distinction permet de poursuivre la construction vers les écrans dimensionnels en conservant l'incidence, la graduation et l'origine de chaque opération.

% BIBITEMS NUEVOS: copiar al entorno thebibliography del artículo.
% Verificación primaria: arXiv 1606.05961v2 (teorema 1.1);
% arXiv 1705.09022v4 (secciones 2--4);
% Conway--Norton, pp. 311 y 314--315, tabla 3.
% \bibitem{chenlamshimakura2016}
% H.-Y. Chen, C. H. Lam y H. Shimakura,
% ``$\mathbb Z_3$-orbifold construction of the Moonshine vertex operator
% algebra and some maximal $3$-local subgroups of the Monster'',
% prepublicación, arXiv:1606.05961 (2016), versión 2 (2017).
% \href{https://arxiv.org/abs/1606.05961}{arXiv:1606.05961}.
% \bibitem{abelamyamada2017}
% T. Abe, C. H. Lam y H. Yamada,
% ``A remark on $\mathbb Z_p$-orbifold constructions of the Moonshine
% vertex operator algebra'', prepublicación, arXiv:1705.09022 (2017).
% \href{https://arxiv.org/abs/1705.09022}{arXiv:1705.09022}.
% \bibitem{kmconwaynorton1979}
% J. H. Conway y S. P. Norton, ``Monstrous Moonshine'',
% \emph{Bulletin of the London Mathematical Society} \textbf{11},
% 308--339 (1979).
% \href{https://doi.org/10.1112/blms/11.3.308}{doi:10.1112/blms/11.3.308}.
